\section{제도와 관련 문헌}

\subsection{한국 경마의 복수 패리뮤추얼 시장}

패리뮤추얼 시장에서는 같은 조합에 걸린 금액이 하나의 풀을 이루고, 공제 후
풀이 적중권에 배분된다. 따라서 게시 배당률은 고정된 북메이커 가격이 아니라
마감 시점의 상대 베팅액을 반영한다. 각 승식은 별도의 풀을 가지므로 공제율,
유동성, 참가자 구성, 배당률 절사 방식의 차이가 교차시장 가격 차이로 남을 수
있다.

\begin{table}[htbp]
\centering
\caption{상위 3위 상태와 주요 승식의 지급 사건}
\label{tab:market-events}
\begin{tabular}{lll}
\toprule
승식 & 지급 사건 & 삼쌍승 상태의 합산 범위 \\
\midrule
단승식 & $i$가 1위 & 고정 $i$, 모든 $j,k$ \\
쌍승식 & $i$가 1위, $j$가 2위 & 고정 $i,j$, 모든 $k$ \\
복승식 & $i,j$가 순서 없이 1·2위 & 고정 $\{i,j\}$, 모든 $k$ \\
삼복승식 & $i,j,k$가 순서 없이 1--3위 & 여섯 순열 \\
삼쌍승식 & $i,j,k$가 정확히 1·2·3위 & 하나의 $(i,j,k)$ \\
\bottomrule
\end{tabular}
\end{table}

연승식과 복연승식은 출전두수에 따라 적중 순위의 범위가 달라질 수 있으므로
주 분석의 단순 주변화 대상에서 제외하고, 공식 승식 규칙을 적용한 별도
강건성 분석으로 다룬다. 이 구분은 편의가 아니라 지급 사건을 정확히 일치시키기
위한 사전 규칙이다.

\subsection{인기마--비인기마 편향과 정보 집계}

\citet{griffith1949}, \citet{weitzman1965}, \citet{ali1977}는 배당률이
시사하는 확률과 실제 승률의 체계적인 차이를 기록했다. 이후 연구는 이 패턴을
위험애호로 읽을 것인지, 낮은 확률의 과대평가로 읽을 것인지 논쟁해 왔다.
\citet{julliensalanie2000}은 위험 아래 선호를 구조적으로 추정했고,
\citet{snowbergwolfers2010}는 위험선호와 확률 오인 설명을 식별하는 검정을
제시했다. \citet{coleman2004}은 여러 국가의 결과를 종합하여 서로 다른 정보와
위험태도를 가진 참가자 집단의 혼합을 한 설명으로 제안했고,
\citet{suhonenetal2018}은 이중이론과 확률가중을 결합하여 패리뮤추얼 자료의
선택을 설명한다.

시장의 차이도 중요하다. \citet{buschehall1988}은 홍콩 시장에서 전형적
편향의 예외를 보고했고, \citet{vaughanwilliams1997}과
\citet{vaughanwilliams1998}은 정보비용, 소비효용, 참가자 이질성이 시장별
편향의 부호를 바꿀 수 있음을 논의한다. \citet{brucejohnson2000}은 병행하는
북메이커와 패리뮤추얼 시장을 비교했고, \citet{sungjohnsonpeirson2012}은
정보가 부족한 참가자의 비중이 가격오류에 미치는 영향을 분석했다.
\citet{shin1991}은 내부정보 거래자가 존재하는 북메이커 모형에서 게시가격이
승률과 체계적으로 벌어질 수 있음을 보였다. 이는 패리뮤추얼 시장을 직접
모형화한 결과는 아니지만, 정보 비대칭이 가격과 기초확률의 괴리를 낳을 수
있는 한 경로를 제시한다.

한국어 문헌도 국내 시장의 서로 다른 측면을 분석했다.
유웅·남준우\citeyearpar{yoonam2015}는 수익률의
평균과 분산에 왜도를 더한 효용함수를 추정하여, 비인기마 선택을 위험애호만으로
해석하기보다 위험회피와 양의 왜도 선호가 함께 나타날 수 있음을 보였다.
같은 저자들\citeyearpar{yoonam2017}은 국내 복승식 자료에서 비선호마 편향과
베팅 전략의 수익률을 분석하여 가격 효율성의 문제를 제기했다.
최혜민 외\citeyearpar{choietal2015}는 경주마·기수·조교사 정보를 사용한 우승마
예측모형을 구축하고, 단승식·복승식·삼복승식의 표본 외 베팅 결과를 비교했다.
이 연구들은 각각 베팅자의 선호, 특정 승식의 가격 효율성, 실현 착순에 대한
외부 예측력을 다룬다.

한국 시장에 대해서는 \citet{jeongetal2019}가 말과 경주의 이질성을 허용하는
승률 모형으로 인기마--비인기마 편향을 검정했다. 본 논문은 이들 연구를 대체하지
않는다. 실현 승률을 추정하거나 개별 승식의 투자수익률을 평가하는 대신, 동일
경주의 복수 승식 가격이 하나의 순위 상태분포에서 나올 수 있는지 검정한다는
점에서 질문과 식별 대상이 다르다.
또한 게시 배당률의 반올림은 인기마 구간의 검정통계에 체계적인 영향을 줄 수
있으므로 \citep{wallsbusche2003}, 원자료의 배당 단위와 상한을 명시적으로
점검한다.
